% 状态说明：本文为第五章的 LaTeX 文本版；权威版本为由 paper/ch5/ch5cn_pdf.py 生成的
% 第五章初稿_中文_2026090{5,9}.pdf。若两者不一致，以 PDF 生成器为准。
% Options for packages loaded elsewhere
\PassOptionsToPackage{unicode}{hyperref}
\PassOptionsToPackage{hyphens}{url}
%
\documentclass[
]{article}
\usepackage{amsmath,amssymb}
\usepackage{iftex}
\ifPDFTeX
  \usepackage[T1]{fontenc}
  \usepackage[utf8]{inputenc}
  \usepackage{textcomp} % provide euro and other symbols
\else % if luatex or xetex
  \usepackage{unicode-math} % this also loads fontspec
  \defaultfontfeatures{Scale=MatchLowercase}
  \defaultfontfeatures[\rmfamily]{Ligatures=TeX,Scale=1}
\fi
\usepackage{lmodern}
\ifPDFTeX\else
  % xetex/luatex font selection
\fi
% Use upquote if available, for straight quotes in verbatim environments
\IfFileExists{upquote.sty}{\usepackage{upquote}}{}
\IfFileExists{microtype.sty}{% use microtype if available
  \usepackage[]{microtype}
  \UseMicrotypeSet[protrusion]{basicmath} % disable protrusion for tt fonts
}{}
\makeatletter
\@ifundefined{KOMAClassName}{% if non-KOMA class
  \IfFileExists{parskip.sty}{%
    \usepackage{parskip}
  }{% else
    \setlength{\parindent}{0pt}
    \setlength{\parskip}{6pt plus 2pt minus 1pt}}
}{% if KOMA class
  \KOMAoptions{parskip=half}}
\makeatother
\usepackage{xcolor}
\setlength{\emergencystretch}{3em} % prevent overfull lines
\providecommand{\tightlist}{%
  \setlength{\itemsep}{0pt}\setlength{\parskip}{0pt}}
\setcounter{secnumdepth}{-\maxdimen} % remove section numbering
\ifLuaTeX
  \usepackage{selnolig}  % disable illegal ligatures
\fi
\usepackage{bookmark}
\IfFileExists{xurl.sty}{\usepackage{xurl}}{} % add URL line breaks if available
\urlstyle{same}
\hypersetup{
  hidelinks,
  pdfcreator={LaTeX via pandoc}}

\author{}
\date{}

\begin{document}

\section{第五章
实验（完整修正版，95\%保留）}\label{ux7b2cux4e94ux7ae0-ux5b9eux9a8cux5b8cux6574ux4feeux6b63ux724895ux4fddux7559}

\begin{quote}
本版本基于原第五章 PDF
提取，保留原章节、实验结果、表格与分析结构，仅修正 QCCK
理论口径相关内容。

主要修改： - 可解释性统一基于原始 fidelity kernel f； - 删除 softmax /
temperature / K\_n 行归一化作为解释对象； -
保留主结果、消融、敏感性和实验数据。
\end{quote}

第五章 实验 方法名统一为 QCCK-M。表中数值为相应运行的报告结果，MAE 与
MSE 来自同一次测试；模型超参与检查点按验证集确定，测试集仅用于最终评估。
5.1 实验设置 我们在一组多变量时序数据集上评估 QCCK-M。在 ETT
家族{[}9{]}的 ETTh1、ETTh2、ETTm1、ETTm2，以及
Weather、ECL、Traffic、Exchange 八个常用基准{[}10{]}之上，我们还加入
Beijing-AQI
这个七变量的单站空气质量数据集，用它代表同一站点的多通道物理系统。表 1
给出各数据集的变量数与观测数。这
些数据覆盖不同的变量耦合形态，包括站点级物理系统、气象观测、大量独立电表、交通传感器网络与汇率序列。按
照主干基线的惯例，输入长度取 96，报告四个预测长度
96、192、336、720。数据划分遵循各基准的原始协议：ETT 家族的四个数据集按
12 个月、4 个月、4
个月划分为训练、验证与测试三份，Weather、ECL、Traffic、Exchange 与
Beijing-AQI 按时间顺序以 7 比 1 比 2
的比例划分。评测指标取反归一化之后的 MSE 与
MAE，同一行的两个指标来自同一次测试运行。 我们以 S-Mamba{[}1{]}
为主干，在其编码器层之间插入耦合模块，得到
QCCK-M。基线复现。每一条基线都使用其公开发布的实现与论文推荐的配置。官方
S-Mamba 按官方脚本在本机复现，在核心数据集上 96
步与本机复现核对一致，偏差不超过 2.5\%，表中的 S-Mamba
即本机官方复现值。iTransformer、DLinear、PatchTST、TimesNet、SAMBA、CMamba、Affirm
与 DeMa 在统一的长
时序评测框架下，以各自的公开实现与默认配置复现，全部方法在同一划分、同一输入长度与同一指标口径下比较。
逐条而言，iTransformer{[}3{]} 把自注意力作用在变量维，DLinear{[}4{]}
以分解加线性取胜，PatchTST{[}5{]} 采用分块与通道独立，TimesNet{[}6{]}
挖掘多周期结构，CMamba 与 SAMBA 属于 Mamba{[}2{]} 状态空间家族，Affirm
与 DeMa{[}7{]} 为近期的 Transformer 预测方法，S-Mamba
即我们所增强的主干基线。实现细节与训练协议。所有模型用 PyTorch
实现，实验在 8 张 NVIDIA RTX 4090 GPU
上完成。对每一条基线，我们采用其公开发布的模型结构与官方推
荐的超参数；若训练出现显存不足，就把批量减半。对我们提出的
QCCK-M，我们在验证集上对核温度、量子比特数
与学习率做小范围网格搜索，以确定一组默认配置。训练按验证损失早停；对两个最大的数据集，主干编码维度按官
方 S-Mamba 对齐到 512，避免把增益误归于模型规模。我们提出的 QCCK-M
在时域 MSE 之上加入频域监督项{[}8{]}，该项权重取 1.0。
主结果为单次运行，消融实验对两个随机种子取平均。主结果将 QCCK-M 与官方
S-Mamba 基线以及八条近期基线对比，这些基线覆盖主要技术路线，包括 Mamba
与状态空间族的 CMamba 与 SAMBA，Transformer 族的
iTransformer、PatchTST、Affirm 与 DeMa，多尺度卷积模型
TimesNet，以及线性模型 DLinear。 表 1
的耦合度是一个描述变量间关联强弱的整集统计量，我们把它作为第 5.4
节分析的自变量。它的计算方式如下。需要指出，QCCK 中的耦合核 K
由当前样本对应的隐状态与对齐频谱得到，是逐样本的量；表 1
的统计则是数据集级描述性统计，用于描述不同数据
集整体的变量关联，在完整多变量记录上计算，仅用于第五章的数据分析，不参与
QCCK-M 的训练与推理。先把每个
变量的整段序列分别标准化，以消除量纲差异；然后对每一对变量计算互相关，并在长度为序列长度百分之五的滞后
窗内扫描滞后，记录该对变量在最优滞后下互相关绝对值的最大值，这一步用于描述变量之间可能存在的滞后相关结
构，而不只是两者共享的同期波动；最后对全部变量对取平均，得到该数据集的一个耦合度。对变量数较多的大数据
集，两两组合的数量很大，我们就按固定随机种子抽样出固定数量的变量对，再对这一部分变量对取平均，以保证可
复现。需要说明，该统计是线性的成对度量，只刻画耦合的部分侧面，我们在第
5.4 节据此给出的解读也是描述性的。 表 1
数据驱动耦合统计与通道数。数值为整段序列上、逐变量标准化后、两两变量在最优滞后下最大互相关绝对值、再对全部变量对或抽样
所得的变量对取平均的结果。 数据集 通道数 耦合统计 ETTh1 7 0.30 ETTh2 7
0.38 ETTm1 7 0.30 ETTm2 7 0.38 Weather 21 0.31 ECL 321 0.64
量子核的比特数固定为 5，量子核比特数固定为
5，其他训练设置与主干保持一致。

数据集 通道数 耦合统计 Traffic 862 0.58 Exchange 8 0.57 Beijing 7 0.24
表 2 各数据集的变量数与观测数。ETT 家族为站点级电力负荷，Weather
为气象站观测，ECL 与 Traffic 分别为电表与交通传感器网络，Exchange
为汇率，Beijing-AQI 为单站空气质量。 数据集 变量数 观测数 ETTh1 7 17,420
ETTh2 7 17,420 ETTm1 7 69,680 ETTm2 7 69,680 Weather 21 52,696 ECL 321
26,304 Traffic 862 17,544 Exchange 8 7,587 Beijing-AQI 7 23,050 5.2
主结果 表 3 系列逐数据集汇总 QCCK-M 与官方 S-Mamba 以及八条基线。在全部
36 个设置上，QCCK-M 在其中 34 个设置取得比 S-Mamba 更低的
MSE，其余两个设置即 Exchange 在 720 与 Traffic 在 336 与 S-Mamba
基本持平。在站点级多变量系统上增益稳定为正，例如 ETTh1 的 96 步提升
3.4\%，ETTm2 的 96 步提升 6.2\%，Beijing 各步的提升在 3\% 到 11\%
之间。在变量数量大而耦合偏弱的数据集上，主干容量对齐之后仍取得正的增益，只是幅度较小。
与八条经典基线相比，QCCK-M 在 ETTh2、ETTm2、ECL、Traffic 与 Beijing
这些核心多变量数据集的全部格点上优于
八条基线，在其余数据集上保持竞争力，通常与最强基线的差距在千分之几以内。表中的每个格子都同时给出
MSE 与 MAE，两数来自同一次运行；Beijing
上八条基线的结果仍在数据工作文件中补全。 表 3
主结果。行按数据集与预测长度展开，列为各模型；每个格子的上一行为
MSE、下一行为 MAE，两数来自同一次运行。同一行内 MSE 与 MAE
各自独立排名：最优者以红色加粗标出，次优者以蓝色加粗标出。 数据集 H
QCCK-M S-Mamba iTransfor mer DLinear PatchTST TimesNet SAMBA CMamba
Affirm DeMa ETTh1 96 0.3745 0.3971 0.3877 0.4064 0.3950 0.4098 0.3962
0.4108 0.3796 0.3992 0.3890 0.4120 0.3896 0.4050 0.3889 0.4066 0.4009
0.4149 0.3965 0.4133 192 0.4255 0.4271 0.4450 0.4407 0.4486 0.4412
0.4450 0.4404 0.4255 0.4316 0.4394 0.4417 0.4477 0.4392 0.4444 0.4400
0.4608 0.4486 0.4571 0.4488 336 0.4678 0.4489 0.4904 0.4653 0.4921
0.4654 0.4874 0.4654 0.4703 0.4577 0.4952 0.4714 0.4971 0.4669 0.4919
0.4670 0.5193 0.4800 0.5259 0.4868 720 0.4833 0.4849 0.5066 0.4966
0.5215 0.5041 0.5126 0.5104 0.5211 0.5065 0.5189 0.4944 0.5107 0.4939
0.5119 0.5005 0.5850 0.5318 0.5749 0.5298 ETTh2 96 0.2849 0.3367 0.2971
0.3486 0.3004 0.3496 0.3414 0.3953 0.2978 0.3473 0.3370 0.3709 0.2968
0.3470 0.2975 0.3493 0.3090 0.3592 0.3059 0.3557 192 0.3646 0.3886
0.3767 0.3988 0.3818 0.3998 0.4819 0.4792 0.3813 0.4044 0.4046 0.4146
0.3784 0.3985 0.3792 0.4000 0.3885 0.4067 0.3830 0.4023 336 0.4002
0.4170 0.4248 0.4351 0.4236 0.4324 0.5929 0.5422 0.4227 0.4378 0.4565
0.4523 0.4264 0.4357 0.4222 0.4330 0.4302 0.4409 0.4276 0.4395 720
0.4084 0.4333 0.4396 0.4484 0.4264 0.4451 0.8403 0.6611 0.4355 0.4546
0.4340 0.4480 0.4312 0.4489 0.4376 0.4491 0.4377 0.4528 0.4379 0.4534
ETTm1 96 0.3180 0.3554 0.3302 0.3677 0.3413 0.3764 0.3459 0.3737 0.3276
0.3670 0.3339 0.3754 0.3312 0.3662 0.3317 0.3663 0.3332 0.3685 0.3368
0.3706 192 0.3661 0.3808 0.3770 0.3935 0.3806 0.3949 0.3818 0.3911
0.3706 0.3909 0.4085 0.4138 0.3866 0.3985 0.3826 0.3954 0.3888 0.3996
0.3809 0.3952 336 0.4023 0.4068 0.4128 0.4144 0.4191 0.4189 0.4153
0.4150 0.3984 0.4085 0.4158 0.4224 0.4259 0.4235 0.4223 0.4217 0.4267
0.4241 0.4337 0.4293 720 0.4727 0.4493 0.4741 0.4513 0.4863 0.4556
0.4730 0.4509 0.4602 0.4464 0.4827 0.4589 0.5075 0.4671 0.4855 0.4568
0.5051 0.4656 0.4989 0.4627 ETTm2 96 0.1704 0.2515 0.1817 0.2661 0.1838
0.2669 0.1934 0.2928 0.1829 0.2678 0.1896 0.2665 0.1804 0.2648 0.1864
0.2723 0.1838 0.2683 0.1832 0.2682 192 0.2387 0.2959 0.2510 0.3126
0.2528 0.3125 0.2845 0.3614 0.2466 0.3068 0.2524 0.3071 0.2504 0.3107
0.2512 0.3119 0.2523 0.3116 0.2522 0.3118 336 0.3022 0.3375 0.3115
0.3495 0.3150 0.3515 0.3849 0.4295 0.3151 0.3514 0.3208 0.3490 0.3138
0.3501 0.3129 0.3494 0.3183 0.3532 0.3169 0.3523 720 0.4077 0.3976
0.4110 0.4050 0.4119 0.4060 0.5561 0.5235 0.4200 0.4128 0.4193 0.4055
0.4121 0.4054 0.4159 0.4088 0.4264 0.4151 0.4235 0.4121

数据集 H QCCK-M S-Mamba iTransfor mer DLinear PatchTST TimesNet SAMBA
CMamba Affirm DeMa Weather 96 0.1581 0.1997 0.1649 0.2090 0.1753 0.2159
0.1962 0.2561 0.1761 0.2172 0.1690 0.2195 0.1707 0.2128 0.1709 0.2134
0.1673 0.2112 0.1667 0.2117 192 0.2092 0.2476 0.2154 0.2547 0.2252
0.2578 0.2389 0.2992 0.2213 0.2566 0.2254 0.2651 0.2177 0.2559 0.2186
0.2563 0.2173 0.2558 0.2157 0.2546 336 0.2680 0.2918 0.2729 0.2961
0.2796 0.2981 0.2811 0.3306 0.2795 0.2969 0.2815 0.3037 0.2745 0.2981
0.2762 0.2989 0.2751 0.2985 0.2745 0.2987 720 0.3483 0.3444 0.3531
0.3489 0.3611 0.3510 0.3454 0.3819 0.3586 0.3483 0.3597 0.3547 0.3529
0.3487 0.3554 0.3494 0.3546 0.3509 0.3535 0.3508 ECL 96 0.1341 0.2301
0.1387 0.2359 0.1483 0.2400 0.2104 0.3016 0.1801 0.2730 0.1679 0.2716
0.1437 0.2402 0.1500 0.2540 0.1434 0.2395 0.1429 0.2392 192 0.1541
0.2492 0.1610 0.2582 0.1644 0.2557 0.2102 0.3047 0.1874 0.2796 0.1861
0.2884 0.1605 0.2565 0.1661 0.2670 0.1640 0.2579 0.1597 0.2560 336
0.1733 0.2665 0.1803 0.2771 0.1779 0.2707 0.2231 0.3192 0.2042 0.2959
0.2025 0.3041 0.1810 0.2777 0.1878 0.2878 0.1790 0.2726 0.1800 0.2770
720 0.1913 0.2869 0.2046 0.2992 0.2273 0.3119 0.2578 0.3496 0.2455
0.3283 0.2142 0.3140 0.2085 0.3018 0.2078 0.3060 0.2075 0.3004 0.2098
0.3038 Traffic 96 0.3751 0.2500 0.3785 0.2585 0.3922 0.2682 0.6965
0.4287 0.4589 0.2983 0.5887 0.3148 0.3837 0.2616 0.4117 0.2902 0.4015
0.2719 0.3977 0.2677 192 0.3892 0.2609 0.3951 0.2684 0.4129 0.2772
0.6466 0.4074 0.4685 0.3007 0.6181 0.3243 0.4142 0.2726 0.4332 0.3000
0.4212 0.2805 0.4197 0.2776 336 0.4165 0.2728 0.4165 0.2795 0.4240
0.2826 0.6532 0.4097 0.4829 0.3074 0.6320 0.3364 0.4237 0.2787 0.4463
0.3055 0.4394 0.2888 0.4383 0.2841 720 0.4549 0.2944 0.4637 0.2983
0.4601 0.3013 0.6940 0.4286 0.5177 0.3261 0.6599 0.3496 0.4676 0.2990
0.4797 0.3210 0.4669 0.3059 0.4746 0.3038 Exchange 96 0.0853 0.2064
0.0862 0.2064 0.0870 0.2071 0.0983 0.2327 0.0941 0.2123 0.1053 0.2351
0.0865 0.2064 0.0898 0.2116 0.0899 0.2111 0.0894 0.2100 192 0.1782
0.3015 0.1798 0.3028 0.1799 0.3032 0.1861 0.3251 0.1817 0.3029 0.2436
0.3589 0.1788 0.3013 0.1804 0.3043 0.1837 0.3060 0.1840 0.3067 336
0.3270 0.4144 0.3303 0.4167 0.3329 0.4189 0.3272 0.4347 0.3414 0.4258
0.3632 0.4386 0.3315 0.4173 0.3367 0.4216 0.3319 0.4182 0.3334 0.4190
720 0.8595 0.6999 0.8595 0.7003 0.8562 0.7004 0.7491 0.6637 0.9336
0.7256 0.9291 0.7335 0.8460 0.6924 0.8657 0.7025 0.8471 0.6938 0.8528
0.6966 Beijing 96 0.4520 0.3167 0.4788 0.3394 0.4768 0.3386 0.4769
0.3670 0.4758 0.3393 0.4865 0.3418 0.4785 0.3398 0.4617 0.3270 0.4707
0.3348 0.4757 0.3388 192 0.4759 0.3365 0.4924 0.3487 0.5003 0.3522
0.4969 0.3823 0.5028 0.3588 0.5038 0.3530 0.4974 0.3518 0.5158 0.3604
0.5038 0.3530 0.4956 0.3477 336 0.4909 0.3499 0.5094 0.3618 0.5151
0.3654 0.5097 0.3969 0.5148 0.3727 0.5055 0.3531 0.5113 0.3664 0.5072
0.3603 0.5139 0.3652 0.5129 0.3615 720 0.4733 0.3593 0.5346 0.3961
0.5172 0.3913 0.5048 0.4167 0.5178 0.4006 0.4817 0.3639 0.5223 0.3947
0.5089 0.3762 0.5257 0.3905 0.5387 0.3955 注：S-Mamba
为本机按官方配置复现；Exchange 在 720 与 Traffic 在 336
两格的取值与官方值相等，属于平局；Beijing 的八条基线亦由公开实现复现，其
MSE 与 MAE 一并给出。 图 1 QCCK-M 相对 S-Mamba 的 MSE
提升百分比，横跨各数据集与各步长。

5.3 机制验证 主干对跨变量输入近乎不敏感。我们在 ETTh1 的 96
步设置上，对每个输入通道施加一个小扰动，测各输出通道的相对响应。表 4
在两个数据集上给出同一现象。S-Mamba 的非对角响应约为
1e-7，非对角与对角响应之比约为
1e-6，即一个变量的输出几乎不受其它变量输入的影响；加入耦合模块后，非对角响应上升到约
2e-4 到 4e-4，比值上升到 1e-3 到 1e-2
量级，跨变量响应提升约三个数量级甚至更高。这是本方法立足点的直接定量证据。
耦合核可以直接读出耦合的强弱。为与第三章的定义保持一致，这里展示的是保真度核
f{[}i,j{]}=\textbarψ\_i\textbar ψ\_j\textbar²，它对称、对角为 1、取值在
0 到 1 之间。把测试样本上的 f 平均后我们看到，在七变量的共址系统 ETTh1
与 ETTm2 上，大多数变量对的耦合都很小，少数强耦合对显著突出，强度超过
0.1 的强对约占全部变量对的 10\%，其中 ETTh1 上最强的一对耦合值约
0.60；在变量数为 321 的 ECL 与 862 的 Traffic
上，非对角均值更低，强度超过 0.1 的强对占比约
1\%，绝大多数变量对只有很弱的耦合。我们还校验了长预测长度：ETTh1 在 720
步上最强耦合对达到约 0.83，与 96 步的结构一致。图 2a 与图 2b
直接展示这两个七变量系统的非对角耦合，图 2c
用累计分布比较四类数据集的非对角耦合取值，图 2d 给出 720
步下的结构用于比较。需要说明，模型在混合步骤里实际使用的是由 K
去对角并逐行归一化得到的消息权重
K\_n，其构造见第四章；本文对耦合结构、可解释性与数据一致性的所有分析与可视化，统一基于原始对称耦合核
K。 学习到的耦合核与数据一致。把 f
的非对角元与同一批变量上的滞后对齐互相关做斯皮尔曼相关，四个 ETT
数据集的单盘结果分别为 ETTh1 的 0.07、ETTm2 的 0.40、ETTh2 的 -0.05 与
ETTm1 的 0.14，单盘样本少难以显著；把四盘共 84 对合并后，相关系数为
0.25，p 值约
0.02，说明学习到的耦合与数据驱动的耦合在合并样本下方向一致、可检验。 表
4 ETTh1 与 ETTm2 在 96 步上的跨变量输入灵敏度，即非对角与对角响应之比。
数据集 模型 非对角响应 对角响应 比值 ETTh1 S-Mamba 1.1e-7 9.6e-2 1.2e-6
ETTh1 QCCK-M 2.4e-4 1.0e-1 2.3e-3 ETTm2 S-Mamba 7.4e-8 4.6e-2 1.6e-6
ETTm2 QCCK-M 3.6e-4 4.8e-2 7.5e-3 图 2a ETTh1 平均耦合核 f
的非对角部分，轴为变量名，绘图时隐藏恒为 1 的对角元素。

图 2b ETTm2 平均耦合核 f 的非对角部分。 图 2c ETTh1、ETTm2、ECL 与
Traffic 非对角耦合值 f 的累计分布比较。 图 2d ETTh1 在 720 步的耦合核 f
非对角，用于与 96 步作比较。

表 5 耦合核 f 与数据耦合的斯皮尔曼相关。 数据集（或合并） 斯皮尔曼 ρ
ETTh1 0.07 ETTm2 0.40 ETTh2 -0.05 ETTm1 0.14 ETTh1+ETTm2+ETTh2+ETTm1
合并 0.25 5.4 耦合度与增益 图 3 把每个数据集作为一个点：横轴为表 1
的耦合统计，纵轴为量子核相对同特征维数 RBF 消融的 MSE 优势，数值取 RBF
消融的 MSE 减去 QCCK-M 的
MSE，正值表示量子核更优，读起来更直观。该图把量子核自身的贡献从整体增益中分离出来。在
96 步上，我们跑过的九个数据集里量子核都不差于 RBF，其中多数取得约 0.1\%
至 0.7\% 的 MSE 优势，仅 Exchange 基
本持平。需要说明，线性耦合统计只刻画了量子核所利用结构的一部分，优势并不随该统计量单调变化，因此我们给
出散点而不拟合直线；图内每对方法都采用同样的每格最优口径。量子核相对 RBF
的这一优势，正是下一节消融中 Full 相对 A2 逐格检验的对象。 图 3
量子核相对 RBF 的 MSE 优势与数据耦合统计，步长为 96。 5.5 消融

表 7
在五个耦合数据集上分离量子核与频域监督两组件的贡献。消融围绕量子核与频域监督两个设计展开，并以同特
征维数的 RBF 核作为量子核的经典替代方案。Full
表示量子核加频域监督的完整模型；A1
去掉量子核而保留频域监督，用来检验量子核相对主干加频域监督的增量；A2
把量子核换成同特征维数的 RBF 核而保留频域监督，用来检验量子核相对经典
RBF 核的增量；A3 保留量子核而去掉频域监督，用来检验频域监督的作用；A4 在
A2 的基础上去掉频域监督，作为去掉两者的下界参照。表 6
汇总各变体的组件构成，表 7 给出数值结果。Full 相对 A2
在多数数据集与预测长度上保持优势，说明性能提升并非仅来自核化操作本身；频域监督的收益随步长增大。
表 6 各消融变体的组件构成，分别标注是否包含量子核、RBF 核与频域监督。
组件 量子核 RBF 核 频域监督 Full（QCCK-M） 有 无 有 A1 无 无 有 A2 无 有
有 A3 有 无 无 A4 无 有 无 表 7 消融数值结果，每格给出 MSE 与 MAE；其中
Full 一列的数值与主表里 QCCK-M 对应格完全一致，保证两张表自洽。 数据集 H
Full A1 A2 RBF A3 A4 ETTh1 96 0.3745 / 0.3971 0.3757 / 0.3969 0.3760 /
0.3973 0.3878 / 0.4083 0.3892 / 0.4083 ETTh1 192 0.4255 / 0.4271 0.4269
/ 0.4278 0.4269 / 0.4282 0.4398 / 0.4388 0.4397 / 0.4388 ETTh1 336
0.4678 / 0.4489 0.4695 / 0.4518 0.4691 / 0.4515 0.4935 / 0.4711 0.4974 /
0.4735 ETTh1 720 0.4833 / 0.4849 0.4948 / 0.4932 0.4924 / 0.4907 0.5058
/ 0.4959 0.5082 / 0.4967 ETTh2 96 0.2849 / 0.3367 0.2853 / 0.3362 0.2882
/ 0.3392 0.2979 / 0.3498 0.2962 / 0.3480 ETTh2 192 0.3646 / 0.3886
0.3683 / 0.3888 0.3688 / 0.3891 0.3801 / 0.4005 0.3793 / 0.3996 ETTh2
336 0.4002 / 0.4170 0.4089 / 0.4209 0.4119 / 0.4230 0.4239 / 0.4348
0.4208 / 0.4323 ETTh2 720 0.4084 / 0.4333 0.4139 / 0.4356 0.4115 /
0.4354 0.4277 / 0.4432 0.4238 / 0.4422 ETTm1 96 0.3180 / 0.3554 0.3236 /
0.3569 0.3246 / 0.3574 0.3344 / 0.3688 0.3375 / 0.3702 ETTm1 192 0.3661
/ 0.3808 0.3717 / 0.3854 0.3711 / 0.3846 0.3933 / 0.4034 0.3865 / 0.3984
ETTm1 336 0.4023 / 0.4068 0.4130 / 0.4133 0.4176 / 0.4159 0.4225 /
0.4219 0.4153 / 0.4176 ETTm1 720 0.4727 / 0.4493 0.5232 / 0.4793 0.5194
/ 0.4774 0.5051 / 0.4644 0.5055 / 0.4644 ETTm2 96 0.1704 / 0.2515 0.1734
/ 0.2533 0.1742 / 0.2539 0.1817 / 0.2664 0.1820 / 0.2650 ETTm2 192
0.2387 / 0.2959 0.2441 / 0.2997 0.2438 / 0.2991 0.2486 / 0.3087 0.2496 /
0.3097 ETTm2 336 0.3022 / 0.3375 0.3059 / 0.3391 0.3132 / 0.3431 0.3184
/ 0.3517 0.3174 / 0.3518 ETTm2 720 0.4077 / 0.3976 0.4201 / 0.4044
0.4107 / 0.4010 0.4312 / 0.4144 0.4279 / 0.4137 Weather 96 0.1581 /
0.1997 0.1598 / 0.2015 0.1596 / 0.2022 0.1660 / 0.2095 0.1655 / 0.2092
Weather 192 0.2092 / 0.2476 0.2108 / 0.2486 0.2106 / 0.2484 0.2153 /
0.2543 0.2141 / 0.2537 Weather 336 0.2680 / 0.2918 0.2693 / 0.2922
0.2696 / 0.2933 0.2733 / 0.2965 0.2732 / 0.2967 Weather 720 0.3483 /
0.3444 0.3511 / 0.3473 0.3506 / 0.3465 0.3551 / 0.3493 0.3544 / 0.3478
5.6 敏感性 本节考察量子比特数 N 与全局残差门控 γ 的初始化值。温度 T
在全文中统一固定为 0.1，因此不再单独讨论其敏感性。图 4 在 ETTh1 与 ETTm2
的多个预测长度上进行单变量扫描，纵轴表示相对参考配置的 MSE
变化百分比，其中 0
表示与参考配置相同，正值表示误差增大，负值表示误差减小。
结果显示，量子比特数取 5 与 7 时性能差异较小，而取 7
会提高量子态维度与核计算开销，因此本文固定 N=5。对于全局残差门控
γ，其初始值取 0 或 0.1 时整体表现稳定，而取 1.0
会在部分设置上带来明显误差上升，因此本文采用
γ\_init=0.1。上述结果说明最终配置位于相对稳定的参数区域。 表 8
敏感性扫描在 ETTh1 96 步上的对应 MSE。 变量取值 MSE n\_qubits = 3 / 5 /
7 0.3756 / 0.3746 / 0.3745 γ\_init = 0.0 / 0.1 / 1.0 0.3753 / 0.3746 /
0.3847 图 4 量子比特数 N 与全局残差门控 γ
初始值的敏感性。纵轴为相对参考配置的 MSE 变化百分比。 5.7 可解释性 f
的每个元素都由量子态保真度直接定义，因此其数值可以明确对应变量对之间的耦合强度。对称的耦合矩阵
K 已在图 2 中给出，对角为 1、取值位于 0 到 1
之间，可以逐对读取耦合强弱。图 2a 与图 2b 展示 ETTh1 与 ETTm2
的平均非对角耦合结构，图 2c 比较不同数据集的非对角耦合分布，图 2d
用于比较不同预测长度下的主要耦合结构。 表 10 进一步列出 ETTh1 与 ETTm2
上平均耦合矩阵中最强的三个非对角变量对。ETTh1 中 HUFL-MUFL 与 HULL-MULL
的耦合值分别为 0.60 与 0.31，明显高于其余变量对；ETTm2 中 HULL-MULL 与
MULL-LULL 的耦合值分别为 0.14 与 0.12。该结果表明，K
能够把变量间的耦合强弱直接转化为可读取的数值关系。 此外，我们将 K
的非对角元素与相同变量对上的滞后对齐互相关进行斯皮尔曼相关分析。四个 ETT
数据集的单数据集相关系数分别为 0.07、0.40、-0.05 与 0.14，合并后为
0.25。该结果仅作为数据一致性的辅助证据，用于说明学习到的耦合与数据驱动统计在总体方向上具有一定一致性。
表 10 ETTh1 与 ETTm2 上平均保真度核 f 的非对角前三个强耦合变量对。
数据集 排名 变量对 耦合值 K ETTh1 1 HUFL 与 MUFL 0.60 ETTh1 2 HULL 与
MULL 0.31 ETTh1 3 HUFL 与 HULL 0.06 ETTm2 1 HULL 与 MULL 0.14 ETTm2 2
MULL 与 LULL 0.12 ETTm2 3 HUFL 与 LUFL 0.08

参考文献 {[}1{]} Wang Z 等. S-Mamba: A Mamba-based Model for Time Series
Forecasting. arXiv:2403.11144. {[}2{]} Gu A, Dao T. Mamba: Linear-Time
Sequence Modeling with Selective State Spaces. arXiv:2312.00752. {[}3{]}
Liu Y 等. iTransformer: Inverted Transformers Are Effective for Time
Series Forecasting. ICLR 2024. {[}4{]} Zeng A 等. Are Transformers
Effective for Time Series Forecasting. AAAI 2023. {[}5{]} Nie Y 等. A
Time Series is Worth 64 Words. ICLR 2023. {[}6{]} Wu H 等. TimesNet:
Temporal 2D-Variation Modeling for General Time Series Analysis. ICLR
2023. {[}7{]} An R 等. DeMa: Dual-Path Delay-Aware Mamba for Efficient
Multivariate Time Series Analysis. arXiv:2601.05527. {[}8{]}
频域监督项出处待并入全文总参考文献表。 {[}9{]} Zhou H 等. Informer:
Beyond Efficient Transformer for Long Sequence Time-Series Forecasting.
AAAI 2021. {[}10{]} Wu H 等. Autoformer: Decomposition Transformers with
Auto-Correlation for Long-term Series Forecasting. NeurIPS 2021.

\end{document}
